Let $H$ be a handlebody.  Let $\dd$ be a disc system for $H$.  Let $D$ be an essential disc in $H$.  The disc $D$ may be very complicated with respect to $\dd$.  In this section we prove Lemma~\ref{subordination-lemma}.  Lemma~\ref{subordination-lemma} lets us construct from $D$ either a disc disjoint from $\dd$ or two ``half-discs'' in $H-\dd$.  In both cases, the constructed objects preserve information about $D$.  Lemma~\ref{subordination-lemma} will be important for later sections.

Let $\dd$ be either a disc system for a handlebody $H$ or a generating system for a free group $F$.  Suppose $\dd$ has $g$ elements.  We use $\dpm$ to denote the set of $2g$ elements consisting of elements of $\dd$ with an orientation.  In the handlebody case, elements of $\dpm$ correspond to the $2g$ sides of discs in $\dd$.  We will usually refer to elements of $\dpm$ simply as discs.  In the free group case, the elements of $\dpm$ correspond to the generators of $F$ and their inverses.  The elements of $\dpm$ generate $F$ as a semigroup.

Define a $\dd$-\emph{interface} to be a partition of $\dpm$ into three subsets labeled black, white, and gray.  We denote these subsets $\db$, $\dw$, and $\dg$.  We require the black and white subsets to be nonempty.  This parallels the definition of interface in Section~\ref{trees}.  We use a single letter, bold and underlined, to denote a $\dd$-interface.  If we want to reference the black and white subsets of $\dpm$ explicitly, we use the notation $\di{\db}{\dw}$.  A \emph{$\dd$-disc} is a $\dd$-interface with $\dg$ empty.  A \emph{$\dd$-half-disc} is a $\dd$-interface with $\dg$ consisting of one element.  The connection to geometric discs and half-discs will become clear during the proof of Lemma~\ref{subordination-lemma}.

Let $H$ be a handlebody of genus $g$.  Let $\dd$ be a disc system for $H$.  Let $\dint{a}$ be a $\dd$-interface.  The $\dd$-interface $\dint{a}$ gives rise to an interface for $\hh$ in the following way.

Consider the tree $\Gamma$ associated to $\dd$.  The ends of $\Gamma$ are canonically in one-to-one correspondence with the ends of $\hh$.  There is an equivariant homotopy $\Phi$ from $\hh$ to $\Gamma$ such that the preimage of the midpoint of any edge equals a lift of the corresponding disc in $\dd$.  The homotopy $\Phi$ identifies the ends of $\hh$ with the ends of $\Gamma$.  

View $\Gamma$ as the Cayley graph of $\pi_{1}(H)$ with generating system $\dd$.  Choose a vertex $v$ of $\Gamma$.  The $\dd$-interface $\dint{a}$ gives colors to the $2g$ edges leaving $v$.  Delete $v$ from $\Gamma$ to form $2g$ connected subsets of $\Gamma$.  Each connected component inherits a color coming from $\dint{a}$.  These components partition the ends of $\Gamma$ and define an interface for $\Gamma$.  This carries over to an interface $\alpha$ for $\hh$.  The interface $\alpha$ is well-defined up to the choice of vertex $v$.  We say $\alpha$ is a \emph{lift} of $\dint{a}$ to $\hh$.  This parallels the notion of lifting a disc in $H$ to an interface for $\hh$.
       
Let $H$ be a handlebody, and let $\hh$ be the universal cover.  Recall from Section~\ref{trees} what it means for a proper curve to cross an interface.  Let $\alpha$ and $\beta$ be two interfaces for $\hh$.  The interface $\alpha$ is \emph{subordinate} to $\beta$ if, for every proper curve $\gamma$, the curve $\gamma$ crosses $\beta$ if $\gamma$ crosses $\alpha$.  In other words, the support of $\alpha$ is contained in the support of $\beta$.  Two interfaces $\alpha$ and $\beta$ are \emph{independent} if no proper curve crosses both.  In other words, they are independent if their supports are disjoint. 

Let $\dd$ be a disc system for $H$.  Let $D$ be a disc in $H$.  Recall what it means to lift a disc to an interface.  A $\dd$-interface $\dint{a}$ is \emph{subordinate} to $D$ if some lift of $\dint{a}$ to $\hh$ is subordinate to a lift of $D$.   Two $\dd$-interfaces $\dint{a}$ and $\dint{b}$ are \emph{independently subordinate}
to $D$ if they can be lifted to independent interfaces that are subordinate to a common lift of $D$.

We can now state and prove Lemma~\ref{subordination-lemma}.

\begin{lemma}
\label{subordination-lemma}
Let $H$ be a handlebody.  Let $\dd$ be a disc system for $H$.  Let $D$ be an essential embedded disc in $H$.  There is either a $\dd$-disc subordinate to $D$ or two $\dd$-half-discs independently subordinate to $D$.
\end{lemma}

\begin{proof}
Let $g$ be the genus of $H$.  Let $\hh$ be the universal cover of $H$.  Let $\Gamma$ be the infinite tree that is homotopic to $\hh$, as constructed above.  The ends of $\Gamma$ are in correspondence with the ends of $\hh$, up to a choice of deck transformation.  Choose such an identification.  Lift $D$ to an interface for $\hh$.  This defines an interface $\alpha$ for $\Gamma$, dividing the ends of $\Gamma$ into  black and white subsets.

Assign the labels black and white to the unoriented edges of $\Gamma$ as follows.  An unoriented edge can be assigned either one, two, or no labels.  Let $(x,y)$ be an oriented edge of $\Gamma$ with origin $x$ and terminal vertex $y$.  Define a \emph{ray} from $x$ to $y$ to be an infinite geodesic path whose initial edge has origin $x$ and terminus $y$.  Every ray in $\Gamma$ determines a unique end of $\Gamma$.

Label the unoriented edge corresponding to $(x,y)$ black (resp. white) if every ray from $x$ to $y$ determines a black (resp. white) end.  If an edge is labeled neither black nor white, label the edge gray.  Since $D$ lifts to a compact
set in $\hh$, the set of gray edges is finite.  Also, the set of gray edges is connected by an easy argument that follows from the definitions.

There are two cases.  Assume first that the set of gray edges is nonempty.  Consider the subgraph of gray edges.  Choose two outermost vertices of this subgraph.  Call the vertices $v$ and $w$.  The $2g$ edges leaving
$v$ are in correspondence with $\dpm$.  Since $v$ is outermost, there is exactly one gray edge leaving $v$.  The other $2g-1$ edges are not all black
and not all white, for otherwise the gray edge would be black or white.  This defines a $\dd$-half-disc that lifts to an interface based at the vertex $v$.  By construction, this interface is subordinate to $\alpha$.  Similarly, we have another $\dd$-half-disc that lifts to an interface based at $w$ and is subordinate to $\alpha$.  Since $v$ and $w$ are outermost vertices of a connected gray subgraph, the two lifts of $\dd$-half-discs are independent.  Thus they are independently subordinate to $D$.

Assume next that there are no gray edges.  Then every edge is colored either black, white, or both.  There are two subcases.  Assume first that some edge $(x,y)$ is both black and white.  Then all rays from $x$ to $y$ define black ends, say, and all rays from $y$ to $x$ define white ends.  Consider the vertex $x$.  Define a $\dd$-disc by setting the oriented edge corresponding to $(x,y)$ black, and the other $2g-1$ edges leaving $x$ white.  This lifts to an interface for $\Gamma$ at $x$ that is equal to $\alpha$, and hence subordinate to $\alpha$.

In the final subcase, each edge has a single label black or white.  Choose a vertex $x$ that touches both a black edge $(x,y)$ and a white edge $(x,z)$.  Then all the rays from $x$ to $y$ are black, and all the rays from $x$ to $z$ are white.  In fact, for any vertex $y$ adjacent to $x$, the rays from $x$ to $y$ are either all black or all white.  Define a $\dd$-disc based at $x$ accordingly.  This defines an interface for $\Gamma$ at $x$ that equals $\alpha$.  This interface is of course subordinate to $\alpha$ and is a lift of a $\dd$-disc.  
\end{proof}
